\section{相关研究与方法选择}
\label{sec:related}
\subsection{传统与生成式匿名化}
本项目把“匿名后能否继续分析表情”作为方法比较的起点。模糊、马赛克和遮挡可以迅速降低面部细节的可见程度，实现成本较低；但它们直接作用于眼睛、眉毛与嘴部时，也会覆盖表情判断所需要的线索。因此，单纯遮住面部并不能同时完成本项目的两个目标。

整脸替换和完整生成提供了另一种思路：通过改变人脸外观削弱身份关联。这类方法的比较重点在于生成内容能否跟随原图的姿态与表情，以及合成范围是否影响头发和背景。对于车内连续画面，面部外区域的保持和相邻帧的连贯性同样需要纳入比较，不能只依据一张生成图的观感选型。

\subsection{关键点与局部编辑}
面部关键点将眉眼、鼻部和嘴部的位置关系表示为几何条件。GANonymization 采用关键点条件引导人脸生成，为保留表情结构提供了可参考的实现\cite{ganonymization}。本项目使用其公开组件，同时通过局部区域合成，将编辑结果限制在人脸内部。

局部残差方法则在灰度人脸基础上学习有限幅度的变化，便于保留亮度、姿态和局部动作。它与几何约束结合后，可以分别控制修改位置和修改强度。项目把这些处理封装为统一匿名化模块，向桌面显示和表情识别提供所需的匿名输出。

\begin{table}[htbp]
\centering
\caption{匿名化备选路线与项目需求的对应关系}
\label{tab:method_compare}
\begin{tabularx}{\linewidth}{@{}lXXX@{}}
\toprule
路线 & 主要处理方式 & 表情信息处理 & 项目选择依据\\
\midrule
模糊与遮挡 & 覆盖或平滑细节 & 关键区域一起被处理 & 用于理解基础匿名方式\\
整脸替换 & 以完整人脸合成 & 需迁移原图动作 & 与局部编辑要求比较\\
关键点条件 & 以几何条件生成 & 由眼眉口结构约束 & 用于匿名化组件\\
局部残差 & 在面部内部编辑 & 保留局部动作线索 & 作为运行链路基础\\
\bottomrule
\end{tabularx}
\end{table}

\subsection{路线选择与评价原则}
综合区域控制、表情保持和计算开销，项目选择局部编辑与关键点约束相结合的路线。人脸检测与对齐负责确定修改位置，软掩膜负责控制合成范围，表情模型负责分析处理后的匿名灰度脸。流程不依赖原图表情预测来选择另一张图片。

实验按照同一任务分别比较匿名化效果、匿名图识别准确率与视频分析速度。前者针对自建身份模型进行项目内对照，后两者使用对应数据集标签和模型处理时间，形成与系统功能直接对应的评价方式。
